\section{Worked contract graph: the root of \texorpdfstring{\(y^3+y=x\)}{y cubed plus y equals x}}
\label{sec:example}

A small example demonstrates the complete architecture without relying on a
large theorem-specific library.  Fix a rational input satisfying
\[
  2^{-a}\leq x\leq2
\]
for a supplied natural apartness parameter \(a\).  The function
\(F(y)=y^3+y\) is strictly increasing on nonnegative rationals, with
\(F(0)=0\) and \(F(1)=2\).  Hence there is a unique real root \(r\in[0,1]\)
of \(F(r)=x\).  We construct finite rational observations of
\[
  r+\frac1r.
\]
The inverse is guarded by the quantitative lower bound on \(x\).

\subsection{Root bisection node}

Start with \([\ell_0,u_0]=[0,1]\).  Given a bracket with
\[
 F(\ell_n)\leq x\leq F(u_n),
\]
let \(m_n=(\ell_n+u_n)/2\).  If \(F(m_n)\leq x\), set
\((\ell_{n+1},u_{n+1})=(m_n,u_n)\); otherwise set
\((\ell_{n+1},u_{n+1})=(\ell_n,m_n)\).

The state uses integer dyadic numerators with a common denominator \(2^n\), so
all operations are primitive recursive and comparisons are exact after
clearing denominators.

\begin{lemma}[Bisection invariant]
For every \(n\),
\[
 0\leq\ell_n\leq u_n\leq1,
 \qquad F(\ell_n)\leq x\leq F(u_n),
 \qquad u_n-\ell_n=2^{-n}.
\]
\end{lemma}

\begin{proof}
The base case is immediate from \(0\leq x\leq2\).  At a step, monotonicity of
\(F\) and the Boolean comparison determine which half still brackets the
root.  Both halves have half the width.  All inequalities are polynomial
inequalities over integers after multiplying by a positive power of two, so
the invariant is quantifier-free.
\end{proof}

A trace package with fuel \(p\) therefore exports a valid root interval of
width \(2^{-p}\).

\subsection{Apartness for the reciprocal}

For \(0\leq r\leq1\),
\[
 x=r^3+r\leq2r,
\]
so
\[
 r\geq x/2\geq2^{-(a+1)}.
\]
This is the crucial guard.  After at least \(a+1\) bisection steps, the dyadic
comparison can certify a positive lower endpoint; more generally the inverse
node requests enough precision that its input interval lies wholly above
zero.

If \([\ell,u]\) is a positive root interval, inversion reverses endpoints:
\[
  r^{-1}\in[1/u,1/\ell].
\]
Suppose the root bracket has width \(2^{-n}\).  If \(n\geq a+2\), its
lower endpoint is at least \(2^{-(a+2)}\), so the reciprocal interval has
width at most
\[
  2^{-n}\,2^{2a+4}.
\]
Consequently a sound reciprocal demand for output precision \(q\) is
\[
 \delta_{\mathrm{inv}}(q,a)=q+2a+4.
\]
The constant is conservative but uniform and is proved by the displayed
rational estimate.

\subsection{Observable node}

Given the root interval \([\ell,u]\) and reciprocal interval
\([1/u,1/\ell]\), output
\[
  B_p=[\ell+1/u,\;u+1/\ell].
\]
It contains \(r+1/r\) by interval arithmetic.  A local rational inequality
bounds its width as a function of the root width and the lower bound
\(2^{-(a+1)}\).  For a requested final precision \(p\), ask the root and reciprocal
outputs at precision \(p+1\).  The backward pass therefore asks the root node at
\[
  p+2a+5.
\]
Each summand interval then has width at most \(2^{-(p+1)}\), so the sum interval
has width at most \(2^{-p}\).

\subsection{Graph and demand flow}

The finite graph is:
\[
 \boxed{x,a}
 \longrightarrow
 \boxed{\text{root bracket}}
 \longrightarrow
 \boxed{\text{guarded reciprocal}}
 \longrightarrow
 \boxed{\text{sum interval}}.
\]
The sum node also reads the root bracket directly.  Backward propagation takes
the maximum of the direct root demand and the stronger demand coming through
the inverse.

For the exact regression instance \(x=1\), \(a=0\), and requested observation
precision \(10\), the supplied Python run computes demands
\[
 d_{\mathrm{observable}}=10,\qquad
 d_{\mathrm{reciprocal}}=11,\qquad
 d_{\mathrm{root}}=15.
\]
It returns the exact root bracket
\[
 \left[\frac{11179}{16384},\frac{22359}{32768}\right]
\]
after fifteen bisections.  The final observable interval has exact width
\[
 \frac{786822173}{8190402920448}<2^{-10}.
\]
These numbers are implementation evidence only; the universal proof is the
bisection invariant plus rational interval estimates.

\subsection{Object-language theorem target}

A closed theorem should quantify over normalized rational codes \(x\), the
apartness exponent \(a\), and requested precision \(p\).  Its conclusion says
that the emitted interval code is well formed, contains every rational
approximation satisfying the root relation at a sufficient finite tolerance,
and has width at most \(2^{-p}\).  A stronger regular-name theorem compares
outputs at two precisions and proves the Cauchy inequality.

A schematic Lean endpoint is:

\begin{lstlisting}[language={}]
def cubicObservableSentence : Fm PTm 0 :=
  all x. all a. all p.
    InputGuard x a -->
    ValidInterval (cubicObservable x a p) p

theorem cubicObservable :
  ([] : List (Fm PTm 0)) |-PRA cubicObservableSentence := by
  exact (PCOG.compile cubicGraph).final
\end{lstlisting}

The actual syntax must expand normalized rational codes, interval validity, and
the polynomial bracket predicate into the repository's \texttt{Fm PTm}
constructors.

\subsection{Proof-node decomposition}

A practical formalization uses the following nodes.

\begin{enumerate}
\item \textbf{Input node}: normalize rational \(x\), prove denominator
positivity, and expose the guard \(2^{-a}\leq x\leq2\).
\item \textbf{Demand node}: compute the general root precision needed for the
final interval width.
\item \textbf{Bisection trace}: maintain bracket and width invariants.
\item \textbf{Lower-bound node}: derive \(r\geq x/2\) at the finite interval
level and prove positivity of the chosen lower endpoint.
\item \textbf{Reciprocal node}: invert interval endpoints and prove
containment.
\item \textbf{Addition node}: add intervals.
\item \textbf{Width node}: prove the final rational width inequality.
\item \textbf{Finalizer}: assemble interval validity and universally close the
parameters.
\end{enumerate}

Only nodes 3 and perhaps 2 require induction.  The rest are exact algebraic
rules and cuts.  This illustrates the ``one induction per recursive component''
principle.

\subsection{What the example teaches}

The example is deliberately modest, but it exhibits four general lessons.

\begin{enumerate}
\item An inverse operation is not justified by a bare nonzero statement; the
input carries a quantitative apartness exponent.
\item Precision flows backward through the graph and can depend on the guard.
\item The root algorithm's correctness is one quantifier-free invariant,
whereas interval arithmetic is acyclic proof composition.
\item The final theorem concerns finite observations and their error, not an
ambient real object.
\end{enumerate}

These are the same structural issues that arise in much larger analytic and
algebraic applications.
